Súčasný pokrok v oblasti umelej inteligencie výrazne mení spôsob, akým ľudia pracujú
s informačnými systémami. Táto transformácia sa prejavuje aj v správe finančných
dokumentov malých a stredných podnikov, kde neustále rastie potreba inteligentných
nástrojov na jednoduchšie vyhľadávanie a interpretáciu údajov.

Veľké jazykové modely (LLM) dnes umožňujú intuitívnu interakciu s dátami prostredníctvom
prirodzeného jazyka. Cieľom tejto diplomovej práce je navrhnúť a implementovať
inteligentného asistenta FinanceGPT integrovaného do existujúcej fakturačnej
aplikácie v Ruby on Rails. Systém dokáže používateľom generovať textové súhrny a vizualizácie
faktúr či výdavkov, pričom zároveň zachováva potrebu používateľskej kontroly pri
kritických výstupoch.

Práca v jednotlivých kapitolách analyzuje možnosti integrácie LLM do webových
aplikácií, navrhuje architektúru systému, dátový model a bezpečnostné opatrenia. Následne
opisuje samotnú implementáciu prototypu a jeho technické overenie na sade evaluačných
požiadaviek. Výsledkom je funkčný nástroj, ktorý zjednodušuje prácu s fakturačnými
údajmi a slúži ako nadstavba pre finančný manažment.